\chapter{Introduzione alla Ricerca Operativa}
\section{Cosa è la Ricerca Operativa?}
La ricerca operativa è la disciplina che si occupa della ricerca nelle operazioni.

La ricerca operativaè, quindi, applicata a problemi di gestione e di coordinamento delle attività nell'ambito di un'organizzazione

\subsection{Significato di "ricerca"}
La parola "ricerca" nel nome evidenzia il fatto che questa disciplina adotta lo stesso approccio usato negli altri campi scientifici.

Infatti, il metodo scientifico è usato, in larga misura, come metodologia di indagine per il problema in esame.

In particolare, il processo inizia con'accurata osservazione e formulazione del problema e la raccolta di tutti i dati relativo.

\section{I modelli nella ricerca operativoa}
Le fasi tipiche di uno studio di ricerca operativa possono essere riassunte in questo modo:
\begin{enumerate}
    \item Definizione del problema e raccolata dei dati
    \item Formulazione del modello matematico
    \item Determinazione delle soluzioni del modello
    \item Test e validazione del modello
    \item Predisposizione di un sistema di supporto all'applicazione del modello
    \item Implementazione del sistema
\end{enumerate}